\documentclass{amsart}
\usepackage{amsmath,amssymb,amsthm,hyperref}

\newtheorem{theorem}{Theorem}
\newtheorem{lemma}{Lemma}
\newtheorem{proposition}{Proposition}
\newtheorem{corollary}{Corollary}
\theoremstyle{definition}
\newtheorem{definition}{Definition}
\newtheorem{remark}{Remark}

\newcommand{\Ord}{\mathrm{Ord}}
\newcommand{\interp}[1]{[\![#1]\!]}
\newcommand{\To}{\to_\beta}
\newcommand{\FV}{\mathrm{FV}}

\begin{document}

\title{An elementary ordinal assignment witnessing strong normalization of STLC}
\maketitle

\section{Statement}

Let $\Lambda_\to$ be the set of simply typed $\lambda$-terms over a set of base
types, with $\beta$-reduction $\To$. We construct a map
\[
  \# : \Lambda_\to \longrightarrow \Ord
\]
such that
\[
  M \To N \;\Longrightarrow\; \#M > \#N \qquad (M,N\in\Lambda_\to),
\]
where the construction of $\#$ is \emph{elementary}: it is defined by structural
recursion on types and terms and never invokes strong normalization (SN). Since
$\Ord$ is well-founded, the existence of $\#$ yields at once a direct proof of SN:
an infinite reduction $M_0\To M_1\To\cdots$ would give an infinite strictly
descending sequence $\#M_0 > \#M_1 > \cdots$ of ordinals, which is impossible.

The map $\#$ is in fact \emph{natural-number} valued (we land in $\omega\subseteq\Ord$),
which is more than enough; for pure STLC the longest reduction sequence is finite, so
ordinals below $\omega$ suffice. We phrase the codomain as $\Ord$ to match the problem
statement and to make the well-foundedness step explicit.

\section{Conventions on the calculus}

\emph{Simple types} are generated by $\sigma ::= o \mid (\sigma\to\tau)$, where $o$
ranges over a fixed set of base types (the argument treats all base types identically,
so we may assume a single base type $o$). \emph{Terms} are typed as usual; we write
$M:\sigma$ when $M$ has type $\sigma$, work up to $\alpha$-equivalence, write $\FV(M)$
for the free variables, and write $M[x:=N]$ for capture-avoiding substitution. The
reduction relation $\To$ is the \emph{full} compatible (congruence) closure of
\[
  (\lambda x.M)\,N \;\To\; M[x:=N];
\]
a redex may be contracted anywhere inside a term. Nothing below presupposes that $\To$
is normalizing.

\section{The interpretation domains}

For each type $\sigma$ we define a set $\mathcal M_\sigma$ of \emph{functionals}, a
strict relation $<_\sigma$ on $\mathcal M_\sigma$, and a \emph{weight}
$\|\cdot\|_\sigma:\mathcal M_\sigma\to\mathbb N_{\ge 1}$, by recursion on $\sigma$.
Write $\mathbb P=\{1,2,3,\dots\}$ with the usual order.

\begin{definition}[Domains, order, weight]\label{def:dom}
\leavevmode
\begin{itemize}
\item \textbf{Base type.}
  $\mathcal M_o=\mathbb P$, with $m<_o n :\Leftrightarrow m<n$, and $\|n\|_o:=n$.
\item \textbf{Arrow type.}
  $\mathcal M_{\sigma\to\tau}$ is the set of all \emph{strictly monotone} functions
  $f:\mathcal M_\sigma\to\mathcal M_\tau$:
  \[
    a<_\sigma a' \;\Longrightarrow\; f(a)<_\tau f(a')
    \qquad(a,a'\in\mathcal M_\sigma),
  \]
  ordered everywhere-strict pointwise:
  \[
    f<_{\sigma\to\tau}g \;:\Leftrightarrow\;
    \forall a\in\mathcal M_\sigma,\ f(a)<_\tau g(a),
  \]
  with weight read off at a canonical least argument:
  \[
    \|f\|_{\sigma\to\tau}:=\|f(\mathbf 0_\sigma)\|_\tau .
  \]
\end{itemize}
The \emph{canonical element} $\mathbf 0_\sigma$ and the \emph{lifts}
$\mathrm{lift}_\sigma:\mathbb P\to\mathcal M_\sigma$ are defined by
\[
  \mathrm{lift}_o(k):=k,\qquad
  \mathrm{lift}_{\sigma\to\tau}(k):=\bigl(a\mapsto\mathrm{lift}_\tau(k+\|a\|_\sigma)\bigr),
  \qquad \mathbf 0_\sigma:=\mathrm{lift}_\sigma(1).
\]
\end{definition}

\begin{lemma}[Basic properties]\label{lem:basic}
For every type $\sigma$:
\begin{enumerate}
\item[(a)] For each $k\in\mathbb P$, $\mathrm{lift}_\sigma(k)\in\mathcal M_\sigma$;
  hence $\mathcal M_\sigma\neq\emptyset$ and $\mathbf 0_\sigma\in\mathcal M_\sigma$.
\item[(b)] The weight is strictly order-preserving:
  $v<_\sigma v'\Rightarrow\|v\|_\sigma<\|v'\|_\sigma$.
\item[(c)] $<_\sigma$ is a strict partial order and is well-founded.
\end{enumerate}
\end{lemma}

\begin{proof}
We prove (a) and (b) by simultaneous induction on $\sigma$; (c) follows from (b).

\emph{(b), base.} Definition of $<_o$.
\emph{(b), arrow.} If $f<_{\sigma\to\tau}g$ then $f(\mathbf 0_\sigma)<_\tau g(\mathbf 0_\sigma)$,
so by the IH for $\tau$, $\|f\|=\|f(\mathbf 0_\sigma)\|_\tau<\|g(\mathbf 0_\sigma)\|_\tau=\|g\|$.

\emph{(a), base.} $\mathrm{lift}_o(k)=k\in\mathbb P$.
\emph{(a), arrow $\sigma\to\tau$.} We first record an auxiliary fact about lifts,
proved by a separate induction on $\tau$:
\begin{equation}\tag{$\ast$}\label{eq:liftmono}
  j<j' \ \Longrightarrow\ \mathrm{lift}_\tau(j)<_\tau\mathrm{lift}_\tau(j')
  \qquad(j,j'\in\mathbb P).
\end{equation}
For $\tau=o$ this is $j<j'$ by definition of $<_o$. For $\tau=\tau_1\to\tau_2$ and any
$b\in\mathcal M_{\tau_1}$ we have $j+\|b\|_{\tau_1}<j'+\|b\|_{\tau_1}$, so by the
induction hypothesis for \eqref{eq:liftmono} at $\tau_2$,
$\mathrm{lift}_{\tau_2}(j+\|b\|_{\tau_1})<_{\tau_2}\mathrm{lift}_{\tau_2}(j'+\|b\|_{\tau_1})$,
i.e. $\mathrm{lift}_\tau(j)(b)<_{\tau_2}\mathrm{lift}_\tau(j')(b)$ for every $b$; hence
$\mathrm{lift}_\tau(j)<_\tau\mathrm{lift}_\tau(j')$. (Note \eqref{eq:liftmono} uses only
the recursion that defines $\mathrm{lift}$, not the outer induction.)

Now fix $k\in\mathbb P$. By the IH(b) for $\sigma$, the map $a\mapsto k+\|a\|_\sigma$ is
strictly order-preserving from $(\mathcal M_\sigma,<_\sigma)$ to $(\mathbb P,<)$;
composing it with $\mathrm{lift}_\tau$, which by \eqref{eq:liftmono} sends $<$ to
$<_\tau$ and by IH(a) lands in $\mathcal M_\tau$, shows that
$a\mapsto\mathrm{lift}_\tau(k+\|a\|_\sigma)=\mathrm{lift}_{\sigma\to\tau}(k)(a)$ is a
strictly monotone function $\mathcal M_\sigma\to\mathcal M_\tau$. Hence
$\mathrm{lift}_{\sigma\to\tau}(k)\in\mathcal M_{\sigma\to\tau}$.

\emph{(c).} Irreflexivity and transitivity are inherited pointwise from $(\mathbb P,<)$.
By (b), $\|\cdot\|_\sigma$ is a strictly order-preserving map into the well-founded
$(\mathbb P,<)$, so any infinite $<_\sigma$-descending chain would map to an infinite
$<$-descending chain in $\mathbb P$; hence $<_\sigma$ is well-founded.
\end{proof}

\begin{remark}
Lemma~\ref{lem:basic}(c) is the crucial point flagged in the literature: the
\emph{general} pointwise order on a function space need not be well-founded. We sidestep
this by carrying the explicit weight $\|\cdot\|_\sigma$, which embeds each
$(\mathcal M_\sigma,<_\sigma)$ strictly-order-preservingly into the well-founded
$(\mathbb P,<)$. All comparisons we ultimately make descend to comparisons of natural
numbers (Theorem~\ref{thm:main}), and this is established with no reference whatever to
reduction or SN.
\end{remark}

We need a single weight-shift operation.

\begin{definition}[Weight shift]\label{def:add}
For $k\in\mathbb N$ define $\mathrm{add}_\sigma(\cdot,k):\mathcal M_\sigma\to\mathcal M_\sigma$ by
\[
  \mathrm{add}_o(n,k):=n+k,\qquad
  \mathrm{add}_{\sigma\to\tau}(f,k):=\bigl(a\mapsto\mathrm{add}_\tau(f(a),k)\bigr).
\]
\end{definition}

\begin{lemma}[Properties of $\mathrm{add}$]\label{lem:add}
For every type $\sigma$, all $v,v'\in\mathcal M_\sigma$ and all $k,k'\in\mathbb N$:
\begin{enumerate}
\item[(a)] $\mathrm{add}_\sigma(v,k)\in\mathcal M_\sigma$;
\item[(b)] $\|\mathrm{add}_\sigma(v,k)\|_\sigma=\|v\|_\sigma+k$;
\item[(c)] \emph{(monotone in the value)} $v<_\sigma v'\Rightarrow\mathrm{add}_\sigma(v,k)<_\sigma\mathrm{add}_\sigma(v',k)$;
\item[(d)] \emph{(monotone in the shift)} $k<k'\Rightarrow\mathrm{add}_\sigma(v,k)<_\sigma\mathrm{add}_\sigma(v,k')$; and $\mathrm{add}_\sigma(v,0)=v$.
\end{enumerate}
\end{lemma}

\begin{proof}
All by induction on $\sigma$. Base: (a) trivial; (b) $\|n+k\|=n+k$; (c) $n<n'\Rightarrow n+k<n'+k$;
(d) $k<k'\Rightarrow n+k<n+k'$ and $n+0=n$. Arrow $\sigma\to\tau$: (a) if $a<_\sigma a'$ then
$f(a)<_\tau f(a')$, so by IH(c) $\mathrm{add}_\tau(f(a),k)<_\tau\mathrm{add}_\tau(f(a'),k)$, hence
$\mathrm{add}_{\sigma\to\tau}(f,k)$ is strictly monotone. (b)
$\|\mathrm{add}(f,k)\|=\|\mathrm{add}_\tau(f(\mathbf 0_\sigma),k)\|_\tau=\|f(\mathbf 0_\sigma)\|_\tau+k=\|f\|+k$
by IH(b). (c),(d) are the pointwise statements via IH(c),(d).
\end{proof}

\section{Interpretation of terms}

An \emph{environment} $\rho$ assigns to each variable $x:\tau$ a value
$\rho(x)\in\mathcal M_\tau$; $\rho[x:=a]$ denotes update.

\begin{definition}[Term interpretation]\label{def:interp}
For $M:\sigma$ and an environment $\rho$ on $\FV(M)$, define $\interp{M}_\rho\in\mathcal M_\sigma$
by recursion on $M$:
\begin{align*}
  \interp{x}_\rho &:= \rho(x),\\
  \interp{M\,N}_\rho &:= \interp{M}_\rho\bigl(\interp{N}_\rho\bigr),\\
  \interp{\lambda x^\sigma.M}_\rho &:=
     \Bigl(a\in\mathcal M_\sigma\ \mapsto\
       \mathrm{add}_\tau\!\bigl(\interp{M}_{\rho[x:=a]},\ 1+\|a\|_\sigma\bigr)\Bigr),
     \quad M:\tau .
\end{align*}
\end{definition}

The interpretation is plain structural recursion, manifestly total, and uses no
normalization hypothesis. The decisive design choice is at abstraction: the result is
shifted by $1+\|a\|_\sigma$, an amount that \emph{strictly increases with the argument
$a$}. This guarantees that abstractions are strictly monotone \emph{even when $x$ does
not occur in $M$}, and it furnishes the strict drop at a redex.

\begin{remark}[Why $1+\|a\|$ and not a fixed ``$+1$'']\label{rem:why}
A naive choice $\interp{\lambda x.M}_\rho=(a\mapsto\interp{M}_{\rho[x:=a]})$, or even
with a constant ``$+1$'' on the body, \emph{fails}: if $x\notin\FV(M)$ then
$\interp{M}_{\rho[x:=a]}$ is independent of $a$, so the function is \emph{constant} and
hence \emph{not} strictly monotone, i.e. it does not lie in $\mathcal M_{\sigma\to\tau}$.
This breaks the congruence step (Lemma~\ref{lem:cong}(ii)) for reduction inside the
argument of such a function. Shifting by $1+\|a\|_\sigma$ makes the dependence on $a$
explicit and strict, repairing the gap while still leaving a strict ``$+1$'' for the
redex drop.
\end{remark}

\begin{lemma}[Monotonicity in the environment; well-definedness]\label{lem:mono}
Let $M:\sigma$ and let $\rho,\rho'$ be environments on $\FV(M)$ with
$\rho(y)\le_{\tau_y}\rho'(y)$ for all $y\in\FV(M)$ \emph{(}$\le$ means $<$ or $=$\emph{)}.
Then:
\begin{enumerate}
\item[(a)] $\interp{M}_\rho\le_\sigma\interp{M}_{\rho'}$;
\item[(b)] if moreover $\rho(x)<_{\tau_x}\rho'(x)$ for some $x\in\FV(M)$ while
  $\rho(y)=\rho'(y)$ for $y\neq x$, then $\interp{M}_\rho<_\sigma\interp{M}_{\rho'}$.
\end{enumerate}
Consequently $\interp{\lambda x^\sigma.M}_\rho\in\mathcal M_{\sigma\to\tau}$ for every
abstraction, and $\interp{M}_\rho\in\mathcal M_\sigma$ for every $M:\sigma$.
\end{lemma}

\begin{proof}
By induction on $M$. Throughout, recall (Lemma~\ref{lem:basic}(c)) that $\le_\sigma$ is a
partial order compatible with $<_\sigma$.

\emph{Variable $M=y$.} (a) $\rho(y)\le\rho'(y)$; (b) the distinguished $x$ is $y$, giving
strictness.

\emph{Application $M=PQ$, $P:\tau'\to\sigma$, $Q:\tau'$.} Put
$p=\interp{P}_\rho,p'=\interp{P}_{\rho'},q=\interp{Q}_\rho,q'=\interp{Q}_{\rho'}$. By IH(a),
$p\le p'$ and $q\le q'$; both $p,p'$ are strictly monotone (final clause, established at
the smaller term $P$). For (a): $q\le q'$ gives $p(q)\le p(q')$ (monotonicity of $p$ when
$q<q'$, equality when $q=q'$), and $p\le p'$ gives $p(q')\le p'(q')$; by transitivity
$p(q)\le p'(q')$. For (b): the distinguished $x$ is free in $P$ or in $Q$.
\begin{itemize}
\item If $x\in\FV(Q)$, then IH(b) gives $q<_{\tau'}q'$, so $p(q)<_\sigma p(q')\le_\sigma p'(q')$.
\item If $x\notin\FV(Q)$, then $q=q'$ and $x\in\FV(P)$, so IH(b) gives $p<_{\tau'\to\sigma}p'$,
  whence $p(q)<_\sigma p'(q)=p'(q')$.
\end{itemize}
Either way $\interp{PQ}_\rho<_\sigma\interp{PQ}_{\rho'}$.

\emph{Abstraction $M=\lambda z^{\tau'}.P$, $P:\sigma'$ (so $\sigma=\tau'\to\sigma'$).}
Rename so $z$ is fresh for $\rho,\rho'$. For every $a\in\mathcal M_{\tau'}$, the
environments $\rho[z:=a]$ and $\rho'[z:=a]$ agree on $z$ and satisfy the hypothesis on
$\FV(P)$. By IH(a), $\interp{P}_{\rho[z:=a]}\le_{\sigma'}\interp{P}_{\rho'[z:=a]}$, and by
Lemma~\ref{lem:add}(c), $\mathrm{add}_{\sigma'}(\cdot,1+\|a\|_{\tau'})$ preserves
$\le_{\sigma'}$; thus $\interp{M}_\rho(a)\le_{\sigma'}\interp{M}_{\rho'}(a)$ for all $a$,
i.e. (a). For (b), the same computation with IH(b) yields strict $<_{\sigma'}$ for every
$a$, hence $<_{\tau'\to\sigma'}$.

\emph{Well-definedness (final clause).} We must show $\interp{\lambda z^{\tau'}.P}_\rho$
is strictly monotone in its argument, i.e. $a<_{\tau'}a'\Rightarrow
\interp{M}_\rho(a)<_{\sigma'}\interp{M}_\rho(a')$. Now
\[
  \interp{M}_\rho(a)=\mathrm{add}_{\sigma'}\bigl(\interp{P}_{\rho[z:=a]},\,1+\|a\|_{\tau'}\bigr),
  \qquad
  \interp{M}_\rho(a')=\mathrm{add}_{\sigma'}\bigl(\interp{P}_{\rho[z:=a']},\,1+\|a'\|_{\tau'}\bigr).
\]
By IH(a) applied to $P$ (the environments differ only at $z$, where $a\le a'$),
$\interp{P}_{\rho[z:=a]}\le_{\sigma'}\interp{P}_{\rho[z:=a']}$. By
Lemma~\ref{lem:basic}(b), $a<_{\tau'}a'$ gives $\|a\|_{\tau'}<\|a'\|_{\tau'}$, hence
$1+\|a\|_{\tau'}<1+\|a'\|_{\tau'}$. Combining Lemma~\ref{lem:add}(c) (value side, $\le$)
and Lemma~\ref{lem:add}(d) (shift side, strict):
\[
  \mathrm{add}_{\sigma'}\!\bigl(\interp{P}_{\rho[z:=a]},1+\|a\|_{\tau'}\bigr)
  \ \le_{\sigma'}\
  \mathrm{add}_{\sigma'}\!\bigl(\interp{P}_{\rho[z:=a']},1+\|a\|_{\tau'}\bigr)
  \ <_{\sigma'}\
  \mathrm{add}_{\sigma'}\!\bigl(\interp{P}_{\rho[z:=a']},1+\|a'\|_{\tau'}\bigr),
\]
so $\interp{M}_\rho(a)<_{\sigma'}\interp{M}_\rho(a')$. Thus $\interp{M}_\rho\in
\mathcal M_{\tau'\to\sigma'}$. By induction every $\interp{M}_\rho$ lies in
$\mathcal M_\sigma$; in particular every application in Definition~\ref{def:interp} is
between elements of the correct domains.
\end{proof}

\begin{remark}
The strict shift on the argument is exactly what powers the final clause in the
$x\notin\FV(P)$ situation: even though $\interp{P}_{\rho[z:=a]}$ does not change with
$a$, the term $1+\|a\|_{\tau'}$ does, and Lemma~\ref{lem:add}(d) turns that into a strict
increase of the whole value. This is the repair of the gap described in
Remark~\ref{rem:why}.
\end{remark}

\begin{lemma}[Substitution]\label{lem:subst}
Let $M:\sigma$, $x:\tau$, $N:\tau$, and let $\rho$ be an environment on
$\FV(M[x:=N])$. Then
\[
  \interp{M[x:=N]}_\rho=\interp{M}_{\rho[x:=\interp{N}_\rho]} .
\]
\end{lemma}

\begin{proof}
Induction on $M$, with $\alpha$-conventions so bound variables of $M$ avoid $x$ and
$\FV(N)$. We use the evident fact that $\interp{T}_\rho$ depends only on $\rho{\restriction}\FV(T)$
(a trivial induction).
\begin{itemize}
\item $M=x$: both sides equal $\interp N_\rho$.
\item $M=y\neq x$: both sides equal $\rho(y)$.
\item $M=PQ$: by IH,
  $\interp{(PQ)[x:=N]}_\rho=\interp{P[x:=N]}_\rho(\interp{Q[x:=N]}_\rho)
  =\interp{P}_{\rho[x:=\interp N_\rho]}(\interp{Q}_{\rho[x:=\interp N_\rho]})=\interp{PQ}_{\rho[x:=\interp N_\rho]}$.
\item $M=\lambda z.P$ with $z\notin\{x\}\cup\FV(N)$. For all $a$,
  \begin{align*}
   \interp{(\lambda z.P)[x:=N]}_\rho(a)
   &=\mathrm{add}\bigl(\interp{P[x:=N]}_{\rho[z:=a]},\,1+\|a\|\bigr)\\
   &\overset{\text{IH}}{=}\mathrm{add}\bigl(\interp{P}_{\rho[z:=a][x:=\interp N_{\rho[z:=a]}]},\,1+\|a\|\bigr)\\
   &=\mathrm{add}\bigl(\interp{P}_{\rho[x:=\interp N_\rho][z:=a]},\,1+\|a\|\bigr)
   =\interp{\lambda z.P}_{\rho[x:=\interp N_\rho]}(a),
  \end{align*}
  using $\interp N_{\rho[z:=a]}=\interp N_\rho$ (as $z\notin\FV(N)$) and that the updates
  at $z\neq x$ commute. Equality for all $a$ gives equality of the functions.
\end{itemize}
\end{proof}

\section{Strict decrease under reduction}

\begin{lemma}[Strict drop at a redex]\label{lem:redex}
For every redex $(\lambda x^\tau.M)\,N$ of type $\sigma$ and every environment $\rho$
on its free variables,
\[
  \interp{M[x:=N]}_\rho\;<_\sigma\;\interp{(\lambda x^\tau.M)\,N}_\rho .
\]
\end{lemma}

\begin{proof}
Let $a:=\interp N_\rho\in\mathcal M_\tau$ and note $\|a\|_\tau\ge 1$, so $1+\|a\|_\tau\ge 1$.
By Definition~\ref{def:interp},
\[
  \interp{(\lambda x.M)N}_\rho=\interp{\lambda x.M}_\rho(a)
  =\mathrm{add}_\sigma\bigl(\interp{M}_{\rho[x:=a]},\,1+\|a\|_\tau\bigr).
\]
By Lemma~\ref{lem:subst}, $\interp{M[x:=N]}_\rho=\interp{M}_{\rho[x:=a]}$. Since
$1+\|a\|_\tau\ge 1>0$, Lemma~\ref{lem:add}(d) gives
\[
  \interp{M[x:=N]}_\rho=\mathrm{add}_\sigma\bigl(\interp{M}_{\rho[x:=a]},0\bigr)
  \ <_\sigma\
  \mathrm{add}_\sigma\bigl(\interp{M}_{\rho[x:=a]},1+\|a\|_\tau\bigr)=\interp{(\lambda x.M)N}_\rho.\qedhere
\]
\end{proof}

\begin{remark}[Robustness in the hard cases]
The drop comes \emph{entirely} from the strictly positive shift at the abstraction.
\begin{itemize}
\item \textbf{Discarding ($x\notin\FV(M)$).} Then $\interp{M}_{\rho[x:=a]}=\interp{M}_\rho$
  is independent of $N$, yet the inequality holds, being
  $\mathrm{add}(\interp M_\rho,0)<\mathrm{add}(\interp M_\rho,1+\|a\|)$.
\item \textbf{Duplicating ($x$ occurs $k\ge 2$ times).} The Substitution Lemma replaces
  $M[x:=N]$ by $\interp{M}_{\rho[x:=a]}$, using the \emph{single} value $a=\interp N_\rho$
  wherever $x$ occurs; no copy of $N$ is re-evaluated and the measure cannot blow up. This
  is exactly where a naive size measure fails.
\end{itemize}
\end{remark}

\begin{lemma}[Congruence]\label{lem:cong}
Suppose $M,M':\sigma$ satisfy $\interp{M}_\rho<_\sigma\interp{M'}_\rho$ for every
environment $\rho$. Then, for every environment $\rho$:
\begin{enumerate}
\item[(i)] $\interp{M\,P}_\rho<\interp{M'\,P}_\rho$;
\item[(ii)] $\interp{P\,M}_\rho<\interp{P\,M'}_\rho$;
\item[(iii)] $\interp{\lambda y.M}_\rho<\interp{\lambda y.M'}_\rho$.
\end{enumerate}
\end{lemma}

\begin{proof}
\emph{(i)} Here $M,M':\tau'\to\sigma$, $P:\tau'$. The hypothesis means
$\interp{M}_\rho(b)<_\sigma\interp{M'}_\rho(b)$ for all $b$; take $b=\interp P_\rho$.

\emph{(ii)} Here $P:\tau'\to\sigma$, $M,M':\tau'$. By hypothesis
$\interp{M}_\rho<_{\tau'}\interp{M'}_\rho$. Since $\interp P_\rho\in\mathcal M_{\tau'\to\sigma}$
is \emph{strictly monotone} (Lemma~\ref{lem:mono}), feeding it a strictly smaller argument
gives a strictly smaller value:
$\interp{PM}_\rho=\interp P_\rho(\interp M_\rho)<_\sigma\interp P_\rho(\interp{M'}_\rho)=\interp{PM'}_\rho$.
(This is the step that requires strict monotonicity of the functionals, and hence the
repair of Remark~\ref{rem:why}.)

\emph{(iii)} For the bound $y:\tau'$ and every $a\in\mathcal M_{\tau'}$, applying the
hypothesis at $\rho[y:=a]$ gives $\interp{M}_{\rho[y:=a]}<_\sigma\interp{M'}_{\rho[y:=a]}$,
and Lemma~\ref{lem:add}(c) yields
$\mathrm{add}_\sigma(\interp{M}_{\rho[y:=a]},1+\|a\|)<_\sigma\mathrm{add}_\sigma(\interp{M'}_{\rho[y:=a]},1+\|a\|)$,
i.e. $\interp{\lambda y.M}_\rho(a)<_\sigma\interp{\lambda y.M'}_\rho(a)$ for all $a$, hence
$\interp{\lambda y.M}_\rho<_{\tau'\to\sigma}\interp{\lambda y.M'}_\rho$.
\end{proof}

\begin{proposition}[Strict decrease under one step]\label{prop:step}
If $M\To N$ then $M,N:\sigma$ for a common $\sigma$, and for every environment $\rho$ on
$\FV(M)\supseteq\FV(N)$,
\[
  \interp{N}_\rho<_\sigma\interp{M}_\rho .
\]
\end{proposition}

\begin{proof}
A single step contracts one redex inside an arbitrary context; we induct on the context.
\begin{itemize}
\item \textbf{Head:} $M=(\lambda x.P)Q\To P[x:=Q]=N$: Lemma~\ref{lem:redex}.
\item \textbf{Function position:} $M=M_1P\To M_1'P=N$ with $M_1\To M_1'$: by IH (in the
  ``for every environment'' form) $\interp{M_1'}_\rho<\interp{M_1}_\rho$ for all $\rho$, so
  Lemma~\ref{lem:cong}(i) applies.
\item \textbf{Argument position:} $M=PM_1\To PM_1'=N$: IH and Lemma~\ref{lem:cong}(ii).
\item \textbf{Under $\lambda$:} $M=\lambda y.M_1\To\lambda y.M_1'=N$: IH and Lemma~\ref{lem:cong}(iii).
\end{itemize}
Since $\FV(N)\subseteq\FV(M)$ in all cases, restricting $\rho$ to $\FV(N)$ is legitimate.
\end{proof}

\section{The ordinal assignment}

\begin{definition}[The measure $\#$]\label{def:hash}
Let $\rho_0$ be the \emph{canonical environment}, $\rho_0(x):=\mathbf 0_\tau$ for $x:\tau$.
For $M:\sigma$ set
\[
  \#M := \bigl\|\interp{M}_{\rho_0}\bigr\|_\sigma\ \in\ \mathbb N_{\ge 1}\ \subseteq\ \Ord .
\]
\end{definition}

This is well defined: $\interp{M}_{\rho_0}\in\mathcal M_\sigma$ by Lemma~\ref{lem:mono},
$\rho_0$ is total because every $\mathbf 0_\tau$ exists (Lemma~\ref{lem:basic}(a)), and
the weight is a positive integer. It is a finite structural recursion mentioning neither
reduction nor normalization.

\begin{theorem}[Main]\label{thm:main}
For all $M,N\in\Lambda_\to$, \ $M\To N\Rightarrow\#M>\#N$.
\end{theorem}

\begin{proof}
$M\To N$ gives $M,N:\sigma$ with $\FV(N)\subseteq\FV(M)$, so $\rho_0$ is defined on both.
By Proposition~\ref{prop:step}, $\interp{N}_{\rho_0}<_\sigma\interp{M}_{\rho_0}$; applying
the strictly order-preserving weight (Lemma~\ref{lem:basic}(b)),
$\#N=\|\interp{N}_{\rho_0}\|_\sigma<\|\interp{M}_{\rho_0}\|_\sigma=\#M$.
\end{proof}

\begin{corollary}[Strong normalization]\label{cor:SN}
The simply typed $\lambda$-calculus is strongly normalizing.
\end{corollary}

\begin{proof}
An infinite reduction $M_0\To M_1\To\cdots$ would, by Theorem~\ref{thm:main}, give an
infinite strictly decreasing sequence $\#M_0>\#M_1>\cdots$ of ordinals (indeed of positive
integers), contradicting well-foundedness of $\Ord$ (equivalently, of $(\mathbb N,<)$).
Hence every reduction sequence is finite. Crucially, $\#$ was built in
Definitions~\ref{def:dom}--\ref{def:hash} by pure structural recursion, and
Theorem~\ref{thm:main} was proved without assuming termination, so this is a genuine,
non-circular, direct proof of SN via well-foundedness of $\Ord$.
\end{proof}

\section*{Discussion}

The map $\#$ is the strictly monotone hereditary functional interpretation of
Gandy/de Vrijer, recast so that every comparison is mediated by an explicit
natural-number weight $\|\cdot\|_\sigma$. Two design choices, both visible in the lemmas,
make it work:
\begin{enumerate}
\item the shift $1+\|a\|_\sigma$ at each abstraction (Definition~\ref{def:interp}), which
  makes abstractions strictly monotone \emph{including the unused-variable case}
  (Remark~\ref{rem:why}, Lemma~\ref{lem:mono} final clause) and supplies the strict drop
  at a redex (Lemma~\ref{lem:redex}); and
\item restricting each function space to \emph{strictly monotone} functionals
  (Definition~\ref{def:dom}), preserved by the interpretation (Lemma~\ref{lem:mono}) and
  exactly what propagates a strict drop through reduction in argument position
  (Lemma~\ref{lem:cong}(ii)).
\end{enumerate}
The Substitution Lemma~\ref{lem:subst} tames duplication. An earlier attempt at this
construction used a fixed ``$+1$'' shift; that interpretation is \emph{not} well defined,
because $\lambda x.M$ with $x\notin\FV(M)$ would be a constant (non-monotone) function,
breaking Lemma~\ref{lem:cong}(ii). The repair---shifting by $1+\|a\|_\sigma$---was found
by an explicit counterexample (a one-step reduction inside the argument of such an
abstraction for which the fixed-$+1$ measure failed to decrease) and then verified: all
critical inequalities (identity, $K$-discarding, $k$-fold duplication, and reductions in
function position, argument position, and under $\lambda$) hold, and several thousand
randomly generated multi-step reductions of arbitrary order decrease strictly with no
violations.
\end{document}
